% !TeX root = ../教材题库_全部条目_转写工作稿.tex
% Source: XJ-B1-C05-TH-P178-Y1217
% 第5章 / 5.3 三角函数的图象与性质 / 教材 p.178、179 / PDF 第185、186页
% OCR未校对：公式与文字均可能有误；未完成数学TeX转写；禁止直接发给学生
% 原题及相邻原解截图（核对用）：../题目原图/XJ-B1-C05-TH-P178-Y1217-01.webp, ../题目原图/XJ-B1-C05-TH-P178-Y1217-02.webp
% 本文件只在转写完成后作为题干使用；OCR 条目尚不是可靠题干。
\begin{bankitem}{p.178 思考（含紧邻说明）}{第5章 / 5.3 三角函数的图象与性质 / 教材 p.178、179 / PDF 第185、186页}{XJ-B1-C05-TH-P178-Y1217}{OCR未校对：公式与文字均可能有误；未完成数学TeX转写；禁止直接发给学生}
\begin{ocrdraft}
[教材p.178 / PDF第185页]
\par
？思考怎样画出余弦函数y一cos的图象？
\par
显然，我们可以仿照上述方法利用余弦线画出y一cos的图象．实际上，由
\par
的图象向左平移《个单位长度，
\par
sinx十一=cos可知，只需把y=sinx
\par
即可得到
\par
y一cos的图象，如图5．3一4．
\par
y=sln盖
\par
y=cos盖
\par
一3此
\par
3涯
\par
4丌
\par
图5．3
\par
一4
\par
[教材p.179 / PDF第186页]
\par
余弦函数的图象称为余弦曲线．
\par
你能由图5，3一得
\end{ocrdraft}
\end{bankitem}
